\subsection{Processo di fornitura}
\label{subsection:processo_fornitura}
Il processo di fornitura, definito dalla norma \glossario{ISO/IEC 12207}, comprende le attività svolte dal \glossario{Fornitore} e inizia con la presentazione di una proposta al Proponente o la stipula di un contratto. Nel nostro caso, il processo riguarda la Proponente Zucchetti S.p.A. e comincia al completamento della fase di accettazione della candidatura.
Il processo prevede la pianificazione e l’organizzazione delle risorse, delle procedure e dei piani necessari per la gestione del progetto, fino alla consegna del sistema, prodotto o servizio software. Questo processo è fondamentale per garantire che il software soddisfi i requisiti del cliente, sia di alta qualità e venga realizzato rispettando tempi e costi concordati. L’obiettivo principale è allineare costantemente le aspettative dell’acquirente con i risultati ottenuti durante l’esecuzione del progetto.
Dunque, per il nostro gruppo, il processo si articola nelle seguenti attività principali:
\begin{itemize}
    \item \textbf{Selezione del capitolato};
    \item \textbf{Candidatura};
    \item \textbf{Pianificazione};
    \item \textbf{Esecuzione e controllo};
    \item \textbf{Revisione e valutazione};
    \item \textbf{Consegna e completamento}.
\end{itemize}
\subsection{Selezione del capitolato}
Il \glossario{Fornitore} va a esaminare i capitolati d'appalto disponibili e in seguito a consultazioni interne raggiunge una decisione unanime sulla candidatura per uno di questi.

\subsection{Candidatura}
Il \glossario{Fornitore} prepara la candidatura al capitolato d'appalto scelto, redigendo i seguenti documenti necessari:
\begin{itemize}
    \item \textbf{Valutazione dei Capitolati}: documento che contiene la valutazione dei vari capitolati d'appalto disponibili, con le motivazioni della scelta;
    \item \textbf{Stima dei Costi e Assunzione Impegni}: documento che contiene un preventivo sulla distribuzione delle ore, la stima dei costi e l'assunzione degli impegni;
    \item \textbf{Lettera di Candidatura}: presentazione del gruppo per il \glossario{Committente}.
\end{itemize}    

\subsubsection{Pianificazione}
Il \glossario{Fornitore} definisce obiettivi, risorse e procedure necessari per l’esecuzione del progetto, identificando i requisiti per la gestione, la misurazione della qualità e lo svolgimento delle attività. Questa fase riguarda principalmente la stesura del Piano di Progetto, che pianifica l’utilizzo delle risorse e documenta i risultati attesi.

\subsubsection{Esecuzione e controllo}
Il \glossario{Fornitore} attua le attività pianificate nel Piano di Progetto, rispettando le norme stabilite. Inoltre, viene effettuato un controllo continuo sullo stato di avanzamento e sulla gestione delle risorse, garantendo la rendicontazione rispetto agli obiettivi prefissati. 

\subsubsection{Revisione e valutazione}
Il \glossario{Fornitore} definisce criteri e procedure per la revisione ed esegue le operazioni in conformità a tali criteri. L'obiettivo è verificare che il progetto soddisfi i requisiti e rispetti gli standard stabiliti.

\subsubsection{Consegna e completamento}
Consegna del progetto, verifica finale e accettazione da parte dell’acquirente.

\subsubsection{Contatti con la Proponente dell’azienda}
La Proponente, Zucchetti S.p.A., fornisce l'indirizzo email del proprio rappresentante per la comunicazione asincrona e per la pianificazione di videochiamate su Google Meet con il \glossario{Fornitore}. 
Le comunicazioni tra Proponente e \glossario{Fornitore} riguardano vari aspetti del progetto, tra cui la raccolta dei requisiti, la raccolta di \glossario{feedback} sui risultati ottenuti e le indicazioni sull'avanzamento del progetto. 
Per ogni colloquio con l'azienda Proponente, ovvero per ogni videochiamata tramite Google Meet, sarà redatto un Verbale Esterno che riassume i punti chiave discussi durante l'incontro.

\subsubsection{Documentazione fornita}
Di seguito viene descritta la documentazione che il gruppo rende disponibile alla Proponente Zucchetti S.p.A. e ai committenti, Prof. Tullio Vardanega e Prof. Riccardo Cardin.

\paragraph{Piano di progetto}
Il Piano di Progetto è un documento redatto dal responsabile del progetto che fornisce una visione dettagliata della gestione e dell’organizzazione del gruppo di lavoro. La sua funzione principale è quella di pianificazione, monitoraggio e controllo delle attività, al fine di garantire il raggiungimento degli obiettivi nei tempi e nei costi stabiliti. Inoltre, il piano include l'analisi e la gestione dei rischi, nonché la pianificazione, il preventivo e il consuntivo di ciascuno \glossario{sprint}, monitorando costantemente l'avanzamento del progetto e le risorse utilizzate.
Il documento è suddiviso nelle seguenti sezioni:
\begin{itemize}
    \item \textbf{Introduzione}: una breve descrizione dello scopo del documento e delle varie sezioni che lo compongono;
    \item \textbf{Analisi dei rischi}: riguarda l'identificazione dei potenziali rischi che potrebbero sorgere durante il corso del progetto, i quali potrebbero causare ritardi od ostacoli nella sua progressione. Vengono inoltre sviluppate strategie di prevenzione per evitare che tali rischi si manifestino, nonché strategie di mitigazione per ridurne l'impatto nel caso in cui si verificassero, al fine di garantire la continuità del progetto;
    \item \textbf{Stima dei costi}: calcolo delle risorse necessarie per il completamento del progetto, che viene aggiornato a ogni \glossario{sprint};
    \item \textbf{Milestone principali}: sezione dedicata alla descrizione delle \glossario{milestone} fondamentali e delle \glossario{baseline} di riferimento del progetto;
    \item \textbf{Primo periodo}: Rappresenta la fase iniziale di avanzamento nel progetto;
    \item \textbf{Sprint}: gli \glossario{sprint} rappresentano periodi di lavoro di durata fissa in cui vengono definiti obiettivi specifici e attività da completare. Durante ogni \glossario{sprint}, il gruppo si impegna a raggiungere tali obiettivi. Inoltre, si stabilisce la data per l'inizio dello \glossario{sprint} successivo.
    La sottosezione si articola nelle seguenti parti:
    \begin{enumerate}
        \item \textbf{Obiettivi}: definisce gli obiettivi specifici dello \glossario{sprint} a livello di \glossario{issue} considerato il pool di issue disponibile;
        \item \textbf{Pianificazione}: presenta, in maniera tabellare, il preventivo dei costi e delle risorse previste per lo \glossario{sprint} indicando anche i tempi di fine preventivati per le varie attività;
        \item \textbf{Consuntivo}: riporta, in maniera tabellare, il resoconto dei costi e delle risorse effettivamente utilizzate;
        \item \textbf{Retrospettiva e valutazione rischi}: riporta la conclusione dello \glossario{sprint} e valuta i rischi che sono emersi, come sono stati mitigati e se ciò non è avvenuto come verranno mitigati;
        \item \textbf{Aggiornamento risorse rimaste}: aggiorna la situazione delle risorse disponibili per il progetto.
    \end{enumerate}
\end{itemize}

\paragraph{Analisi dei requisiti}
L’analisi dei requisiti è un documento redatto dall’analista che fornisce una visione chiara delle richieste e delle aspettative dell’azienda Proponente. Va a includere un elenco dettagliato delle funzionalità da sviluppare e implementare, nonché i casi d’uso che definiscono le interazioni tra il sistema e l’utente.
Il documento è suddiviso nelle seguenti sezioni:
\begin{enumerate}
    \item \textbf{Introduzione}: Una breve descrizione dello scopo del documento e dei riferimenti usati per la sua stesura;
    \item \textbf{Descrizione del prodotto}: Descrive gli obiettivi del prodotto, le sue funzioni principali e le caratteristiche;
    \item \textbf{Requisiti utente}: Elenco completo dei requisiti utente trovati, suddivisi in obbligatori e opzionali;
    \item \textbf{Use case}: Indica tutti i casi d’uso individuati dal gruppo durante l'analisi;
    \item \textbf{Requisiti software}: Elenco completo dei requisiti software del prodotto, suddivisi in funzionali e non funzionali(di qualità e di vincolo).
\end{enumerate}

\paragraph{Piano di Qualifica}
Il Piano di Qualifica è il documento che tratta della specifica degli obiettivi di qualità di prodotto e processo. 
Definisce un insieme di indici di valutazione e validazione del progetto che verranno usati durante lo svolgimento dello stesso per garantire qualità. 
Il documento è suddiviso nelle seguenti sezioni: 
\begin{enumerate}
    \item \textbf{Introduzione}: Una breve descrizione dello scopo del documento e delle varie sezioni che lo compongono;
    \item \textbf{Obiettivi di qualità}: Questa sezione presenta i valori accettabili e gli ambiti per le \glossario{metriche} definite dal team. Viene divisa per metriche e ogni sotto sezione è dotata di una tabella che descrive le metriche con: ID della metrica, nome, valore tollerabile e valore ambito;
    \item \textbf{Cruscotto di valutazione della qualità}: Vengono visualizzati i grafici relativi ai valori delle \glossario{metriche} e una loro analisi. Il cruscotto di qualità viene
    aggiornato in maniera periodica per allineare le valutazioni agli \glossario{sprint} più recenti.
\end{enumerate}

\paragraph{Glossario}
Il Glossario è un elenco dettagliato e organizzato di tutti termini, acronimi e definizioni utilizzati nella documentazione. 
L'obiettivo principale di questo documento è fornire una comprensione chiara dei concetti e dei termini specifici impiegati nel progetto, garantendo una comunicazione coerente e precisa tra tutti i membri del gruppo e con gli \glossario{Stakeholder}. 
In questo modo si facilita il lavoro collaborativo e si assicura un allineamento efficace su linguaggio e significati durante l'intero \glossario{ciclo di vita} del progetto.

\paragraph{Lettera di presentazione}
La Lettera di Presentazione è il documento con cui il gruppo comunica la propria candidatura alle revisioni di avanzamento \glossario{Requirements and Technology Baseline} e \glossario{Product Baseline}. Nel primo caso essa include informazioni sui \glossario{repository} di documentazione e codice sorgente, il riferimento al \glossario{Proof of Concept} e un aggiornamento sugli impegni con la stima del preventivo "a finire". 

\paragraph{Specifica Tecnica}
La Specifica Tecnica è un documento che tratta la progettazione e l'architettura del software del prodotto. Va a descrivere le scelte progettuali e le tecnologie usate per lo sviluppo di quest'ultimo. 
Viene suddiviso nelle seguenti sezioni:
\begin{itemize}
    \item \textbf{Introduzione}: Una breve descrizione dello scopo del documento e dei riferimenti usati per la stesura;
    \item \textbf{Tecnologie utilizzate}: Descrive le tecnologie utilizzate per lo sviluppo del prodotto;
    \item \textbf{Architettura}: Descrive l'architettura del prodotto, con particolare attenzione ai pattern architetturali usati;
    \item \textbf{Progettazione ad alto livello}: Descrive ad alto livello le varie componenti del sistema e le sue funzionalità;
    \item \textbf{Progettazione del front-end di dettaglio} descrive in dettaglio i componenti del front-end;
    \item \textbf{Progettazione del back-end di dettaglio} descrive in dettaglio i componenti del back-end. Include anche i diagrammi delle classi e i design pattern utilizzati.
\end{itemize}

\paragraph{Manuale Utente}
Il Manuale Utente è un documento che fornisce tutte le informazioni necessarie all'utente finale per utilizzare il prodotto software. Il manuale è redatto in modo chiaro e comprensibile, con l'obiettivo di guidare l'utente passo dopo passo nell'utilizzo del software. 
Il documento è suddiviso nelle seguenti sezioni:
\begin{itemize}
    \item \textbf{Introduzione}: Una breve descrizione dello scopo del documento e dei riferimenti usati per la stesura;
    \item \textbf{Requisiti minimi per l'uso}: Specifica i requisiti minimi necessari per l'utilizzo del software;
    \item \textbf{Installazione}: Guida l'utente nell'installazione del software;
    \item \textbf{Avvio}: Spiega come avviare il software;
    \item \textbf{Utilizzo}: Spiega all'utente come utilizzare il software, con esempi pratici e illustrazioni;
\end{itemize}

\subsubsection{Strumenti}
Qui vengono elencati gli strumenti usati, che vengono poi approfonditi nella sezione \hyperref[sec:strumenti]{Strumenti}.
\begin{itemize}
    \item LaTeX;
    \item Git;
    \item Zoom;
    \item Meet;
    \item Google Fogli;
    \item Google Documenti.
\end{itemize}
